
We have presented an NK landscape framework for studying voting methods as collective optimizers. The main findings are:

\begin{enumerate}
  \item \textbf{Complexity-dependent method ordering.} Under direct democracy,
    method performance undergoes sharp phase transitions: total score $\to$
    scoring ($p\!=\!0.35$) $\to$ Borda $\to$ STAR as complexity increases. No
    method universally dominates.

  \item \textbf{Borda's robustness.} Borda count is simultaneously the most
    efficient and most equitable method across a wide complexity range. A sweep
    over the generalized scoring family reveals that $p=0.35$ (a flatter
    ordinal scoring rule) dominates plurality in the low-to-moderate regime while matching
    Borda's favorable variance properties.

  \item \textbf{Voting rules as utility transforms.} The model frames voting
    methods as transformations of a common utility matrix, suggesting that formal
    connections between transform properties and landscape-dependent optimality
    may be possible.

  \item \textbf{Representation as information bottleneck.} Under representative
    democracy, the regime structure is reshaped by candidate selfishness and
    identity voting, with simpler methods gaining advantage as representation
    degrades.
\end{enumerate}

These results suggest that debates over voting reform are implicitly debates about the complexity of the policy landscape. Future work might explore strategic voting, develop formal connections between utility transform properties and landscape-dependent optimality, introduce fitness-biased candidate selection to model incumbency advantage, and extend the framework with a third, socially transmitted ``belief'' portion of the bitstring so that opinion dynamics on networks can interact with collective optimization.
